% =============================================================================
% Project : THE_PHILOSOPHY_OF_VIABILITY_LIBRARY_PROJECT
% Book    : Book I — Ontology of Viability
% Book ID : VS.B01
% Part    : Part I — Identity and Manifestation
% Chapter : Chapter 08 — Comparative Philosophical Examination
% Chapter ID : VS.B01.P01.C08
%
% File :
%   4_ASSETS/tables/
%   VS.B01.P01.C08.S03.SS07.P003.TL01.heraclitean_interpretive_positions.tex
%
% Section / Chapter / Asset:
%   Table TL01 — Heraclitean Interpretive Positions
%
% Structural Identifier:
%   VS.B01.P01.C08.S03.SS07.P003.TL01
%
% Asset Type : TL
%
% Purpose:
%   To compare major interpretive approaches to Heraclitus,
%   their textual basis, explanatory contribution, and limits.
%
% =============================================================================

\begin{table}[htbp]
\centering
\small

\caption{Principal Interpretive Approaches to Heraclitus}
\label{tab:VS.B01.P01.C08.S03.SS07.P003.TL01}

\begin{tabularx}{\textwidth}{
    >{\raggedright\arraybackslash}p{2.55cm}
    >{\raggedright\arraybackslash}p{2.7cm}
    Y
    Y
}
\toprule
\textbf{Interpretive approach}
&
\textbf{Principal textual emphasis}
&
\textbf{Explanatory contribution}
&
\textbf{Main interpretive limitation}
\\
\midrule

Flux-centred reading
&
River fragments; elemental transformation; B36 and B91
&
Highlights change, turnover, and the instability of particular
conditions.
&
May overgeneralize from selected fragments and treat every form
of continuity as illusory.
\\

Material-monist reading
&
Fire, cosmic exchange, and elemental transformations; B30, B31, B90
&
Explains the importance of fire in the cosmological fragments
and connects Heraclitus with earlier material-principle theories.
&
May reduce transformation to changing appearances of one
underlying material substance.
\\

Order- and measure-centred reading
&
Logos, measure, and the common order; B1, B2, B30, B50
&
Explains why change is associated with order rather than
unrestricted disorder.
&
May impose a fully systematic metaphysical principle where
the surviving evidence remains compressed and ambiguous.
\\

Unity-of-opposites reading
&
Contextual and reciprocal relations; B10, B51, B61, B88
&
Shows how contrasting conditions can be related without
being simply identical.
&
May treat distinct forms of opposition as though they all
instantiate one universal logical structure.
\\

Process-oriented reading
&
Transformation, reciprocal change, and continuity; B12, B31, B88
&
Highlights the possibility that continuity and transformation
may be internally related.
&
May translate Heraclitus too quickly into later process
metaphysics or a general theory of persistence.
\\

\bottomrule
\end{tabularx}

\caption*{\footnotesize
The approaches are interpretive orientations rather than mutually
exclusive schools. Their evidential support and explanatory scope
remain matters of scholarly dispute. The table does not present
any one interpretation as definitively established.
}

\end{table}

% =============================================================================
% END OF TABLE TL01
% =============================================================================